\chapter{主成分分析与特征表示}
\chaptergoal{解释相关特征与低维表示的关系；理解方差最大方向、主成分和解释方差比例；使用 sklearn 的 PCA 分析 Wine 和 Digits；比较降维前后模型的表现；正确处理 PCA 与交叉验证的数据边界。}

科学数据常包含许多互相关联的字段。多波段测光、频带能量与图像像素都可能存在冗余。如果用少量变量概括主要变化，绘图和后续建模可能更方便。主成分分析（principal component analysis, PCA）提供一种经典的线性降维方法，但“用较少维度保留更多方差”与“保留最适合科学目标的信息”并不是同一个判断。
\section{相关特征与低维表示}

一条记录用几个数来描述，就称它在当前表示中有几个\textbf{维度}。葡萄酒的13个输入对应13维，手写数字的一幅图展平成64个数则对应64维。这里的维度是数据表示的长度，不是物理空间有多少维；同一个对象可以根据选用特征具有不同的数据维数。

\textbf{降维}是把原来较多的输入数值变换成较少的新数值，保留对某种目标有用的结构。\textbf{特征选择}是挑出原来的部分列；PCA的\textbf{特征变换}则把原列进行线性组合，产生新的坐标。例如新变量$z=(x_1+x_2)/\sqrt2$混合了两列，不等于只留下$x_1$。主成分分析选择哪些组合，要由数据的变化方向决定。

以 Wine 的13个特征为例，如果某些化学测量指标随样本整体变化趋势相近，可能存在相关结构。高维数据的每条记录可视为 $\mathbb R^p$ 中的一个点；许多点的变化可能主要集中在少数方向上。降维试图找到这些主要方向，再用新的坐标描述样本。

对于天文测光，某些星等或颜色可能相关；对于声呐，邻近频带的能量可能具有连续结构。但相关结构不保证来自单一物理原因，也可能受共同测量误差、仪器响应或样本选择影响。PCA 首先是一种数据表示方法，而不是物理成分识别仪器。

\subsection{为什么压缩不只是删除几列}
设两个测量量近似满足 $x_2\approx x_1$，散点主要沿对角线展开。删除任意一列可以保留一部分变化，但沿对角线的平均方向可能更有效地概括两个测量共同的信息。PCA构造的是新坐标，不一定从原始列中挑选某一列。新坐标仍需记录其与原变量的关系，才能用于新数据。

相关性也不是重复的充分条件。两个高度相关的波段可能在差值中包含重要颜色信息；差值的变化很小，却恰好区分类别。降维必须先明确目标是压缩、可视化、去噪还是改善预测。PCA直接优化的方差保留与平方重构，不是所有这些目标的同义词。

\textbf{秩}描述数据矩阵中独立线性方向的数量。若第二列严格等于第一列的两倍，两列只提供一个独立线性方向；若存在噪声，通常不再严格低秩，但某些方向变化很小，可以考虑近似低秩。这个概念帮助区分精确代数冗余与允许信息损失的压缩。

\subsection{先中心化，再讨论方向}
PCA研究相对于平均位置的变化。若不中心化，离原点很远的均值会影响第一方向，使它更多反映整体偏移而不是样本之间的变化。中心化使用训练均值，不需要标签；新样本也减同一个均值。否则每批新数据会拥有不同原点，投影坐标不能直接交给同一个后续模型。

对均值为 $\boldsymbol\mu$ 的输入，记 $\mathbf x_c=\mathbf x-\boldsymbol\mu$。中心化不是标准化：前者只移动原点，后者还改变各坐标的长度单位。不同选择会定义不同的“最大变化”问题，应根据字段单位及科学目标决定。

\section{投影与方差的几何解释}

\textbf{中心化}是每列减去该列训练均值，使新的坐标原点位于样本均值处。\textbf{投影}可理解为从一个点向某条方向线作垂线，用垂足在线上的位置作为新的坐标。若方向向量$\mathbf v$长度为1，中心化点$\mathbf x_c$在该方向上的坐标为$z=\mathbf x_c^{\mathsf T}\mathbf v$。这里$z$是一个数，$\mathbf v$是一组方向系数，两者不是同一个对象。

取四个人工二维点$(-2,-2),(-1,-1),(1,1),(2,2)$，两列均值都是0，所以已中心化。选择方向$\mathbf v=(1,1)/\sqrt2$，四个投影坐标分别为$-2\sqrt2,-\sqrt2,\sqrt2,2\sqrt2$。投影的样本方差为
\[
\frac{(-2\sqrt2)^2+(-\sqrt2)^2+(\sqrt2)^2+(2\sqrt2)^2}{4-1}=\frac{20}{3}.
\]
与它垂直的方向$(1,-1)/\sqrt2$上，四个投影全为0，方差为0。全部变化都沿第一条线，因此本例一个坐标就足够表达原来的四个点。

\fig{v2_pca_projection.pdf}{人工二维点向主方向投影。另加的查询点$(1,0)$用于展示有信息损失的投影，不参与确定主方向。}

要重构原坐标，把$z$乘回$\mathbf v$再加均值即可。例如$(2,2)$的$z=2\sqrt2$，可准确重构。另一个不在线上的点$(1,0)$投影后$z=1/\sqrt2$，重构成$(0.5,0.5)$，丢失了垂直于线的分量，平方误差为$0.5$。真实数据通常不像四点例子完全共线，所以压缩会丢失部分信息。

在二维平面中，让一组点向一条指定直线作正交投影。选择不同的投影方向，投影坐标的分布宽度会变化。PCA 寻找使投影后样本方差尽可能大的方向作为第一主成分；第二主成分与第一主成分正交，在其余方向中保留尽可能多的方差，以此类推。

对中心化后的数据矩阵 $X_c$，第一主成分可写为单位方向 $\mathbf v_1$，使 $\|X_c\mathbf v_1\|^2$ 最大。这里强调“投影”和“方差”即可，求解特征向量的完整推导可以选读。主成分方向的符号并不唯一，把某个主成分整体乘以$-1$不会改变所保留的信息。

\subsection{从给定方向走向由数据求方向}
前面的共线例子可以直观看出方向，但一般数据不会恰好落在一条直线上。现在用一组不完全共线的人工点，完整计算协方差矩阵，再得到主要变化方向。

\worked{四个二维点的协方差与特征方向}
取 $(2,1),(1,2),(-1,-2),(-2,-1)$。两列均值都为0。以样本协方差的 $n-1=3$ 为分母，有
\[
S=\frac1{3}X^\mathsf TX=
\begin{pmatrix}10/3&8/3\\8/3&10/3\end{pmatrix}.
\]
对角项来自各列平方和10，非对角项来自乘积和8。方向 $\mathbf v_1=(1,1)^\mathsf T/\sqrt2$ 满足 $S\mathbf v_1=6\mathbf v_1$，正交方向 $\mathbf v_2=(1,-1)^\mathsf T/\sqrt2$ 满足 $S\mathbf v_2=(2/3)\mathbf v_2$。因此特征值为6和 $2/3$，主要变化沿对角线。

四个第一坐标组成向量
\[
\mathbf z=(3/\sqrt2,\;3/\sqrt2,\;-3/\sqrt2,\;-3/\sqrt2)^\top.
\]
只保留这一坐标后，前两点均重构为$(1.5,1.5)$，后两点均重构为$(-1.5,-1.5)$。每条原记录到重构点的平方距离为0.5，总和为2。若按每条记录的二维平方距离取平均，误差为0.5；若对全部8个坐标元素取均方误差，则为0.25。两种归一化都可以，但不能混称为同一个数。

\fig{v3_ch13_geometry.pdf}{不完全共线的四个人工点、由协方差求出的第一方向及秩1重构。方向由数据计算，不是预先指定一条便于绘图的直线。}

\subsection{为什么最大方差方向是最大特征值方向}
单位方向 $\mathbf v$ 上的投影方差为 $\mathbf v^\mathsf TS\mathbf v$。设协方差具有一组正交单位特征向量，特征值按 $\lambda_1\geq\lambda_2\geq\cdots\geq0$ 排列，把 $\mathbf v$ 展开成 $\sum_ja_j\mathbf v_j$，则 $\sum_ja_j^2=1$，并有
\[
\mathbf v^\mathsf TS\mathbf v=\sum_j\lambda_j a_j^2\leq\lambda_1.
\]
选择 $\mathbf v=\mathbf v_1$ 可以达到上界。这个短推导把“寻找最宽方向”与线性代数中的特征向量联系起来。第二方向在与第一方向正交的条件下重复同样选择，依次得到后续方向。

若最大的两个特征值相等，最大方差方向可能不唯一。即使特征值不同，每个方向仍可以整体乘负号；投影得分也同时变号，重构不变。因此不同软件或不同运行出现整条主轴翻转，不一定意味着错误。需要比较子空间和重构，而不是机械比较每个方向符号。

\section{主成分和解释方差}

\textbf{主成分方向}是PCA从数据中求得的新坐标轴方向；沿该方向的投影值叫主成分得分或新坐标。描述各原特征怎样参与线性组合的系数，有时称为载荷；不同资料对载荷是否另外乘以方差尺度存在约定差异，本书画方向系数时明确指单位方向向量。

\textbf{解释方差}是投影坐标在样本中的方差，不是“模型能解释多少物理因果”。假设某数据的两个主方向方差为4和1，总和为5，第一主成分的解释方差比例是$4/5=80\%$，两个一起为100\%。这个百分数描述输入数据的变化量，不是80\%的分类正确率。

设第 $j$ 个主成分对应的解释方差为 $\lambda_j$，解释方差比例为
\begin{equation}
q_j=\frac{\lambda_j}{\sum_{k=1}^{p}\lambda_k}.
\end{equation}
前 $m$ 个主成分的累计比例为 $\sum_{j=1}^m q_j$。解释方差越高，表示保留的原始数据总体变化越多，但并不自动意味着预测目标的信息越多。低方差方向也可能包含区分类别的重要信号，尤其当不同类别在总体数值变化很小的方向上分开时。

选择维度时应考虑可视化需要、重构误差和下游任务表现，而不应机械要求“累计解释方差达到95\%才算正确”。若后续是监督分类，维度选择必须在开发数据的验证流程中进行。

\subsection{方向、坐标和解释方差分别是什么}
第一主成分方向是一组作用于原始特征的系数 $\mathbf v_1$；第 $i$ 个样本在该方向上的得分是 $z_{i1}=\mathbf x_{ci}^\mathsf T\mathbf v_1$。系数是模型学到的方向，得分是每条记录变换后的坐标，二者维度不同。某些文献把方向系数称为载荷，另一些把它乘以特征值平方根后称为载荷，使用时应注明约定。

解释方差比例为 $r_j=\lambda_j/\sum_k\lambda_k$，累计保留比例为前 $q$ 个 $r_j$之和。在四点例子中，第一方向占 $6/(6+2/3)=0.9$。这表示保留90\%的总样本方差，不表示90\%记录重构完全正确，也不表示分类准确率90\%。每个对象都可能损失一部分信息。

对中心化矩阵 $X_c$作奇异值分解 $X_c=U\Sigma V^\mathsf T$，$V$的列给出同样的方向，非零特征值为 $s_j^2/(n-1)$。这里介绍SVD是为了连接软件求解方式和前面的几何量，不要求学生推导数值算法。PCA通常可以通过这种分解求得，而不是先猜方向再用梯度下降逐个调整。

\subsection{方差最大与平方重构最小的关系}
沿单位方向的正交投影满足勾股分解：一个中心化样本的总平方长度等于保留坐标平方加垂直残差平方。对所有样本相加，数据总平方长度固定，因此最大化保留投影平方和，就是最小化重构残差平方和。在选取 $q$ 个正交方向时，最优总残差为
\[
\|X_c-\widehat X_c\|_F^2=(n-1)\sum_{j>q}\lambda_j.
\]
$\|\cdot\|_F^2$表示矩阵所有元素平方之和。四点例子丢弃特征值 $2/3$，得到 $3\times2/3=2$，与逐点计算一致。这是一条可以交叉核对的关系，而不是额外规定的评价指标。

这个最优性以线性、正交投影和平方重构为前提。它不说明PCA是所有非线性压缩、所有噪声模型或所有科学任务的最优表示。

\worked{保留近乎全部方差，也能丢掉全部分类信息}
取四点 $(-10,-0.1),(-10,0.1),(10,-0.1),(10,0.1)$，标签由第二坐标正负决定。第一坐标方差远大于第二，原始尺度PCA保留第一方向可保留约99.99\%方差，但在每个第一坐标上都同时有两类，无法依据这个坐标分开。第二坐标虽变化很小，却给出完美的人工类别规则。

这个例子没有否定PCA；它说明压缩目标与预测目标不同。若为了分类选择保留维数，应将PCA放进完整预测验证流程，而不是仅凭累计解释方差超过95\%就宣布选择完成。

\section{标准化与 PCA 流程}
PCA 对变量尺度敏感。若一个字段的数值范围远大于其他字段，不经处理时它可能主导总方差。教学中对 Wine 使用标准化再做 PCA，帮助学生认识尺度处理的作用。但若各物理量的原始方差本身具有科学意义，标准化是否合理仍需基于问题决定。

\begin{lstlisting}[caption={Wine 的 PCA 可视化准备}]
from sklearn.datasets import load_wine
from sklearn.preprocessing import StandardScaler
from sklearn.decomposition import PCA

X, y = load_wine(return_X_y=True)
from sklearn.model_selection import train_test_split
X, X_hold, y, y_hold = train_test_split(
    X, y, test_size=0.2, stratify=y, random_state=42
)
X_std = StandardScaler().fit_transform(X)
coords = PCA(n_components=2).fit_transform(X_std)
print(coords.shape)
\end{lstlisting}

上述拟合只使用开发数据，仅供\textbf{无监督表示的探索性绘图}。如果要用 PCA 的结果作为监督分类器输入，并报告泛化分数，就必须把 StandardScaler、PCA 和分类器放在一个 \texttt{Pipeline} 中，在每个训练折内独立拟合。

\fig{ch13_wine_pca.pdf}{Wine 标准化后前两个主成分的投影，使用原始类别作事后着色以便解释。主成分本身没有使用类别标签学习。}

\subsection{尺度决定PCA认为哪些方向重要}
假设一列长度用厘米、另一列质量用千克，直接比较原始方差包含单位选择。如果把长度改为毫米，方差乘100，主方向可能明显变化，但实验对象没有改变。标准化后PCA比较的是各变量相对自身训练波动的变化，通常更适合量纲不同的表格特征。

但标准化也不是无条件正确。若各列同单位、噪声水平相近而方差差异正是信号强弱，统一到单位方差可能放大低信噪比列。对于光谱像素，逐波长标准化会改变原来的强弱与连续谱结构；对于有误差估计的测量，还可以考虑噪声加权，但那已是另一种明确的优化问题。书中葡萄酒默认标准化，是针对不同指标量级的教学选择。

\subsection{在每折内同时学习缩放与方向}
PCA即使不使用标签，也会学习数据分布。若在全部样本上先拟合主轴，再对变换结果交叉验证，验证记录已经影响了表示。对未来未见样本的评价，应在每个训练折重新拟合中心、尺度和方向，并对该折验证记录只作变换。

\begin{lstlisting}[caption={把方向选择放进每一折，而不是先对全数据降维}]
from sklearn.pipeline import make_pipeline
from sklearn.preprocessing import StandardScaler
from sklearn.decomposition import PCA
from sklearn.linear_model import LogisticRegression
from sklearn.model_selection import cross_val_score
model = make_pipeline(
    StandardScaler(), PCA(n_components=3),
    LogisticRegression(max_iter=2000))
scores = cross_val_score(model, X_dev, y_dev, cv=cv)
\end{lstlisting}

保留维数3也是一个待选择配置。若比较1、2、3、5等值，必须承认验证结果参与了选择。图中为了探索而在全部开发数据上拟合的PCA，可以用于描述开发分布，但不能与严格折内拟合的评价混为同一个训练过程。

\subsection{白化改变坐标的距离含义}
PCA通常只旋转再截断坐标，保留方向上的尺度差异。\textbf{白化}进一步按主成分标准差缩放，使保留坐标具有相近方差。它可能帮助某些距离模型，也可能放大本来很小且受噪声主导的方向。报告中必须说明是否使用白化，不能把PCA、标准化和白化三个操作当作完全相同的“归一化”。

\section{葡萄酒测量数据的投影与分析}
课堂实验先检查原始各特征的尺度，再计算 PCA 解释方差。绘制前两个主成分的散点图，并根据原有类别作不同标记。图中可能出现较明显的类别分离，但这不等于证明无监督 PCA 已经完成分类；类别信息只用于可视化解释。分类性能需要训练分类器并在独立验证上评价。

\begin{lstlisting}[caption={将 PCA 纳入监督分类流程，避免折间泄漏}]
from sklearn.pipeline import make_pipeline
from sklearn.model_selection import cross_val_score, StratifiedKFold
from sklearn.linear_model import LogisticRegression
X_dev, y_dev = X, y  # development data from the previous cell
cv = StratifiedKFold(5, shuffle=True, random_state=7)
for components in [2, 5, 10, "no_pca"]:
    steps = [StandardScaler()]
    if components != "no_pca":
        steps.append(PCA(n_components=components))
    steps.append(LogisticRegression(max_iter=2000))
    pipe = make_pipeline(*steps)
    score = cross_val_score(pipe, X_dev, y_dev, cv=cv).mean()
    print(components, score)
\end{lstlisting}

这个示例在同一开发部分通过交叉验证观察维度的影响，保留部分没有参与选择。维度选择后的最高验证分数不能直接当作最终独立测试结果。

其中\texttt{"no\_pca"}分支完全不加入PCA，用作真正的不降维参照。\texttt{PCA(n\_components=None)}一般仍会中心化并旋转到所有主方向，不能把“保留全部成分”误说成“未执行PCA”。
上述代码已经先留出测试数据，只在开发集内选择主成分数。

\subsection{实际解释方差与预测表现要分别报告}
在142条葡萄酒开发记录上拟合标准化和PCA，前两方向累计方差约0.5506，前五个约0.8008，前八个约0.9209。这些是开发数据的表示描述，不使用类别来确定主轴。画图时可以用类别标记作事后观察，但要说明标签不是PCA的训练目标。

\fig{v3_ch13_components.pdf}{葡萄酒开发数据的累计解释方差。图只回答表示保留多少方差；分类效果通过另一个折内拟合实验评价。}

在相同五折中，将不同PCA维数与逻辑回归串联，得到准确率：1维约0.8599，2维约0.9719，3维约0.9581，5维约0.9719，8维和13维均约0.9788。保留方差随维数单调增加，分类准确率却不是单调增加。这与前面的目标区分一致。

全13维时，相当于在标准化空间作完整正交变换。对适当的线性模型，完整旋转不损失输入信息；数值差异仍可能来自求解容差和具体惩罚设置。真正的信息压缩发生在丢弃方向时。与直接标准化逻辑回归比较时，应保持训练折和正则化选择范围一致。

\subsection{怎样解释一个主方向}
可以查看方向系数，找出哪些原特征同向变化、哪些相反，再回到字段字典解释这些测量组合。系数大不自动等于特征对分类重要；主方向是无监督变化方向，不是类别判别方向。相关字段也可能共同反映仪器或样本选择，不一定对应一种化学过程。

若两个主方向接近简并，轻微数据变化可以使它们在同一二维子空间中旋转，而累计解释方差和重构几乎不变。此时解释整个子空间可能比逐个命名“第一主成分的物理意义”更稳妥。

\section{手写数字图像重构与高维扩展}
Digits 每张图片是 $8\times8$ 的灰度像素，展平后有64个数值特征。把主成分数从64降到10或20，可减少存储表示维度，再通过 PCA 的逆变换近似恢复图像。比较原图与重构图能够直观理解“保留主要变化”与“丢失细节”的含义。

\fig{ch13_digits_recon.pdf}{Digits 的原始像素与使用有限主成分重构的示例。重构显示维度压缩会丢失部分图像细节。}

在图像、光谱或传感器时间序列上，PCA 既可用于压缩，也可用于查找主要变化模式。但这些模式常混合物理信号、观测条件和噪声，需要结合专业背景解释。

\subsection{用看得见的误差检查压缩}
手写数字数据的一条记录是8乘8灰度图，展平后64列对应固定像素位置。只在训练图像上拟合PCA，再对验证图像变换和逆变换，可以观察不同维数保留了什么。重构值可能不是整数，也可能超出原始灰度范围；这属于线性重构的结果，显示时裁剪与计算误差前裁剪是不同操作，应明确区分。

本轮固定开发内部划分，使用原像素尺度，验证集每像素MSE在2、8、20、40个成分时分别约为13.4181、6.2947、2.0785、0.2526。它们是所有验证图像和64个像素的平均平方误差，不能与前面按每个二维点的平方距离直接比较数值。

维数增加使本次重构误差下降，但某个数字的关键笔画仍可能在较低维时被抹去。应同时查看代表性重构、逐图误差与下游分类，而不是只展示最容易辨认的几幅图。如果依据验证图选择维数，它又成为一个开发选择，最终评价需保留独立记录。

\subsection{从图像到科学序列}
光谱也可组织成“对象为行、波长位置为列”的矩阵，但在使用PCA前必须确保各列对应可比的波长，处理缺失、分辨率、归一化及测量误差。直接拼接不同波长采样的数组，会让列之间失去共同含义。时序也类似，需要区分固定时刻、相位对齐或提取后的特征表示。

主成分重构可以作为压缩或异常残差工具，但残差大既可能是新物理，也可能是校准、坏像素或对齐问题。下一章研究无监督分组，第15章再讨论异常候选的解释，不能把“PCA没解释好”直接当成新类型证明。

\section{专业特征与后续模型比较}
天文方向可用 SDSS 的测光特征观察降维结构，物理方向可用 Sonar 的60个频带能量检查主成分和分类表现。如果原数据的主要变化由某个工况或仪器响应主导，PCA 找到的主要方向也可能首先反映该因素，而不是目标类别。可通过将主成分得分与已知测量条件作图辅助解释，但相关不能直接证明因果。

\subsection{一个有控制的专业迁移方案}
以声呐能量特征为例，先建立标准化加逻辑回归或SVM基线，再把PCA作为一个候选步骤加入，比较固定开发折内的准确率、目标类召回率与计算成本。每次训练折都重新学习PCA。若用完整数据的标签图挑选“最好看的两个成分”，再在同一数据上报告分类，图的选择已使用额外信息，应明确其探索性质。

对天文测光特征，注意星等和颜色并非独立测量：多个颜色可能共享同一波段误差。PCA可能集中展示相关误差或整体亮度方向，而不是消除观测偏差。应查看各方向与亮度、信噪比及质量字段的关系，同时保持用于评价的对象划分。

\subsection{本章应形成的三项不同判断}
第一，线性代数判断：方向和投影是否算对，重构误差是否满足相应关系。第二，统计判断：保留多少变化，方向在抽样扰动下是否稳定。第三，任务判断：这种表示是否有利于预测、可视化或科学分析。三者可以支持彼此，但不能用其中一个数字代替另外两个判断。

\chapterpractice{必做：Wine 的标准化、解释方差与二维投影；在同一开发/验证方案下比较不降维与不同 PCA 维度的分类器。提高：Digits 原图与重构图；专业迁移：Sonar 或 SDSS 的特征表示。}
\chapterexercise
\begin{enumerate}
\item 为什么 PCA 第一个主成分通常优先保留较大的数据方差？
\item 主成分解释方差达到90\%是否保证分类准确率没有下降？举出可能反例。
\item 为什么用整表训练 PCA 再做交叉验证可能导致评价失真？
\item 对两个存在显著不同物理单位的变量，标准化前后 PCA 的方向可能怎样变化？
\item PCA 中某主成分和物理量高度相关，能否直接断言主成分对应单一物理机制？
\end{enumerate}

\section*{深化练习与方法检验}
本组题强调解释、计算与验证；D编号在本章内独立计数。计算题应写出所用定义和中间步骤，实验题保留实际输出，包括不支持预期的结果。
\begin{enumerate}[label=D\arabic*.]
\item 对正文非共线四点 $(2,1),(1,2),(-1,-2),(-2,-1)$，计算样本协方差矩阵及两个特征值。用一维重构核对丢失的总平方距离。
\item 一个第一主成分在两个软件中符号相反。怎样检查它们是否表示同一子空间？若两个特征值几乎相等，只比较各轴的逐元素差异是否合理？
\item 构造两维数据：第一维有很大的无关噪声方差，第二维数值变化很小但决定类别。说明为什么按解释方差保留一个主成分可能损失分类能力，以及怎样验证。
\item 若训练集有30条记录、100个特征，中心化后的数据矩阵秩最多是多少？这怎样约束非零主成分数，为什么不能把100个输入自动理解为100个独立变化方向？
\end{enumerate}
